\cotefile{93}
\scanpage{135}
la quelle extension s'identifie à $\mathrm{P}^{1}_{X/S}(\omega)$
par sa structure d'extension naturelle de
$\omega$ par $\underline{\omega}\otimes\deleted{\unsure{\underline{\omega}}}\Omega^1_{X/S}\simeq\underline{\omega}\otimes\underline{\omega}=\underline{\omega}^{2}$) s'identifie
à l'op. diff. de degré $1$ (rel. à $S$) universel
sur $\underline{\omega}$
\[
C\vert\underline{\omega}=d^{1}_{\underline{\omega},X/S}:
\underline{\omega}\longrightarrow{\unsure{\mathscr P|\mathrm P|\mathcal P}}^{1}_{X/S}(\underline{\omega}).
\]

L'indétermination, \inserted{donc,} en cette construction,
au lieu d'être dans $\Gamma(\underline{\omega}\otimes\mathscr J)$
\deleted{\unsure{est}} (opérant sur $\mathscr J$ par crochet, à valeurs
dans $\mathscr J\otimes\underline{\omega}$) est \deleted{dans} \inserted{une section de}
$\underline{\omega}\otimes\mathrm{Fil}_0\simeq\underline{\omega}\otimes\underline{\omega}\simeq
\underline{\omega}^{\otimes 2}$ -- on retombe sur ses pattes !

\smallskip\centerline{\rule{1.2cm}{0.35pt}}\smallskip

Ceci ``rappelé'', on voit donc que si $S=\operatorname{Spec}\mathbb C$,
\deleted{\unsure{la|donc}} $X$ courbe algébrique lisse $/\mathbb C$, alors
la \deleted{\unsure{argument}}\inserted{construction} \uline{transcendante} \deleted{\unsure{appliquée}}\inserted{indiquée}
plus haut \deleted{\unsure{nous fournit de|nous fournit des}}\inserted{s'interprète comme la} construction
d'une connexion projective (canonique)
sur $X/\mathbb C$, i.e. d'une section canonique
du torseur ${\unsure{\mathscr C|C|\mathcal C}}_{X/\mathbb C}$ \inserted{sous $\underline{\omega}_{X/\mathbb C}^{\otimes 2}=\underline{\omega}^{2}$} des connexions
projectives sur $X$. \inserted{On se demande quand} cette connexion
projective \deleted{\unsure{est-elle|est}} de plus ``régulière à
l'infini'' i.e. la connexion correspondante sur
$\mathbb P_X\simeq\mathbb P(\mathscr J/\mathscr J^{3})$ ($\mathscr J$ Idéal
d'augmentation de $\mathscr P^\infty_{X/\mathbb C}$) \deleted{étant}\inserted{est} régulière
à l'infini. Sauf erreur, cela doit signifier

% margin note written sideways
\margin{par un argument analogue on voit que sur une courbe elliptique privée d'un pts, la connex. projective canonique n'est pas régulière, \unsure{Voir|Vois} (\unsure{autrement} elle ne \unsure{dépendrait} pas du dit trou, par raison d'\unsure{homogénéité} -- \unsure{trop rapide|très rapide} !)}
% slanted
\margin{Ce n'est pas le cas sur ${\unsure{E|\mathbb E}}_{\mathbb C}^{1}$ à $\unsure{0|1}$ \unsure{trou}, \unsure{ni} sur $\mathbb P^1_{\mathbb C}$ privée de 3 pts, \unsure{enfin} (car sur $\mathbb P^1_{\mathbb C}$ \unsure{privé de} $\nu\leq3$ pts, il n'y a au \unsure{plus} qu'une seule conn. proj. régulière (à l'$\infty$ !))}
% archive sheet number: 134
